\section{Tests unitarios}
\label{sec:unit}

Cada caso se ha aislado en un programa de ensamblador propio, generado con el
\sig{script.py} provisto en \texttt{asm\_test/} (excepto \texttt{tb\_abort}, que se
distribuye con la ROM precompilada \sig{memoriaRAM\_I\_abort.vhd} para situar la
rutina de servicio en la palabra 96 de la ROM, exigida por el vector de \sig{Data\_abort}).
La tabla~\ref{tab:fsm-cobertura} resume qu\'e estados y transiciones de la FSM ejercita
cada test, y d\'onde se observa la verificaci\'on en el cronograma.

\renewcommand{\arraystretch}{1.18}
\begin{table}[htbp]
\centering
\footnotesize
\begin{tabularx}{\textwidth}{@{}llXl@{}}
\toprule
\textbf{Test} & \textbf{Estados ejercitados} & \textbf{Casos verificados (hit/miss, set, v\'ia)} & \textbf{Cronograma} \\
\midrule
\texttt{tb\_write\_hit}
  & \estado{Inicio}, \estado{Arbitro}, \estado{ADDR}, \estado{Fetch}
  & 1 read miss s0v0, 5 hits (3 read + 2 write); \sig{cont\_w} se actualiza
    s\'olo en write hit; bloque queda dirty.
  & \dato{TBD ciclo $t_1$} \\
\texttt{tb\_write\_miss}
  & + \estado{WriteAround}
  & 1 write miss (write-around) que no sube bloque y no toca \sig{cont\_w};
    1 read miss posterior trae el dato actualizado.
  & \dato{TBD ciclo $t_1$} \\
\texttt{tb\_dirty\_replacement}
  & + \estado{CopyBack}
  & Llena s0v0+s0v1, ensucia v0, fuerza miss s0 con reemplazo dirty
    (\sig{cont\_cb}++); luego reemplazo limpio en el mismo set.
  & \dato{TBD ciclo $t_1$} \\
\texttt{tb\_scratch}
  & + \estado{Scratch}
  & 2 sw + 2 lw a MD Scratch (single word, no cacheable). Verifica que MC
    intacta tras los accesos.
  & \dato{TBD ciclo $t_1$} \\
\texttt{tb\_io}
  & \estado{Inicio} (s\'olo)
  & sw \sig{Output\_Reg}, lw \sig{Input\_Reg}, lw \sig{Output\_Reg}, sw
    \sig{ACK\_Reg}. Sin actividad en bus de memoria.
  & \dato{TBD ciclo $t_1$} \\
\texttt{tb\_abort}
  & FSM error: \estado{No\_error}$\to$\estado{memory\_error} y vuelta
  & lw desalineado dispara \sig{Mem\_Error}; lectura de \sig{Addr\_Error\_Reg}
    en RT\_Abort lo limpia. Variantes (b)(c) documentadas.
  & \dato{TBD ciclo $t_1$} \\
\bottomrule
\end{tabularx}
\caption{Cobertura de la FSM por unit test. Tras simular cada caso, sustituir
\dato{TBD} por la marca temporal del cronograma (en ns) donde se observa el
evento principal.}
\label{tab:fsm-cobertura}
\end{table}
\renewcommand{\arraystretch}{1.0}

\subsection{Resultados esperados de los contadores}

\begin{table}[htbp]
\centering
\begin{tabular}{lcccc}
\toprule
\textbf{Test} & \sig{cont\_m} & \sig{cont\_w} & \sig{cont\_r} & \sig{cont\_cb} \\
\midrule
\texttt{tb\_write\_hit}          & 1  & 2  & 6  & 0 \\
\texttt{tb\_write\_miss}         & 3  & 1  & 4  & 0 \\
\texttt{tb\_dirty\_replacement}  & 4  & 1  & 7  & 1 \\
\texttt{tb\_scratch}             & 2  & 0  & 4  & 0 \\
\texttt{tb\_io}                  & 1  & 0  & 3  & 0 \\
\texttt{tb\_abort}               & 1  & 0  & 1+\dato{$n$} & 0 \\
\bottomrule
\end{tabular}
\caption{Valores finales esperados de los contadores (verificar al final de
cada simulaci\'on). En \texttt{tb\_abort} \sig{cont\_r} crece cada vez que
RT\_Abort vuelve a leer \sig{Addr\_Error\_Reg}, por eso queda en funci\'on
de cu\'antas iteraciones del bucle de error se simulen.}
\label{tab:contadores-esperados}
\end{table}

\subsection{Comentarios por caso}

\paragraph{\texttt{tb\_write\_hit}.} Sirve para visualizar el camino corto
\estado{Inicio}$\to$\estado{Arbitro}$\to$\estado{ADDR}$\to$\estado{Fetch}$\to$\estado{Inicio}
una vez (el read miss inicial) y posteriormente cinco transiciones autom\'aticas
\estado{Inicio}$\to$\estado{Inicio} (todos los hits). Las dos escrituras hit
deben dejar el bit \sig{dirty} del bloque a `1' (visible en
\sig{Via\_2026\_CB.dirty\_bits}).

\paragraph{\texttt{tb\_write\_miss}.} Demuestra la pol\'itica write-around: el sw a
\texttt{0x40} pasa por \estado{Arbitro}$\to$\estado{ADDR}$\to$\estado{WriteAround}
y vuelve a \estado{Inicio} sin actualizar la cach\'e. El siguiente lw a la misma
direcci\'on s\'i sube el bloque (read miss limpio), demostrando que la escritura
fue efectiva en MD.

\paragraph{\texttt{tb\_dirty\_replacement}.} El paso (3) entra a
\estado{CopyBack} (con \sig{send\_dirty}=`1' y \sig{MC\_bus\_Write}=`1' ya en
\estado{ADDR}), transfiere las 4 palabras del bloque sucio, vuelve a \estado{ADDR}
para emitir la fase de direcci\'on del fetch y completa con \estado{Fetch}. La
pareja \estado{CopyBack}+\estado{Fetch} es el camino m\'as costoso de la m\'aquina.

\paragraph{\texttt{tb\_scratch}.} \estado{Scratch} se distingue de \estado{WriteAround}
en que sirve tanto lecturas como escrituras (\sig{MC\_bus\_Read}=\sig{RE},
\sig{MC\_bus\_Write}=\sig{WE}). Se observa que el \sig{addr\_non\_cacheable} hace
saltar desde \estado{Inicio} a \estado{Arbitro} sin incrementar \sig{cont\_m}.

\paragraph{\texttt{tb\_io}.} El decodificador interno de \sig{IO\_MD\_subsystem}
intercepta los accesos a \texttt{0x000070xx} antes de la MC, por lo que las
se\~nales \sig{RE\_MEM} y \sig{WE\_MEM} se quedan a `0' en estos accesos: la FSM
permanece en \estado{Inicio} y los contadores no cambian.

\paragraph{\texttt{tb\_abort}.} El acceso desalineado en \estado{Inicio} dispara
inmediatamente la transici\'on de la FSM secundaria a \estado{memory\_error}
(\sig{load\_addr\_error}=`1', \sig{Mem\_Error}=`1'). El procesador salta al
vector de \sig{Data\_abort} y la rutina situada en la palabra 96 lee
\sig{Addr\_Error\_Reg}, lo que en \estado{Inicio} dispara la transici\'on
\estado{memory\_error}$\to$\estado{No\_error}.
